\begin{abstract}
\noindent
We decompose an H\&M-anchored Hassan--Rosen bimetric benchmark across DESI DR2 BAO, Pantheon+ supernovae and a compressed Planck~2018\,+\,ACT~DR6 CMB likelihood. At the adopted anchor, BAO is indistinguishable from \LCDM{} ($\Dchi=0.179$), whereas SNe retain $\Dchi=2.94$ with the full STAT+SYS covariance, dominated by the quartic-closure density mismatch rather than by direct late-time shape. The late-time background follows \LCDM{} with effective density $\cos^2\theta\,\Om$, with relative residual at most $3.4\times10^{-5}$ in $E^2$ for $z\le2.33$; this common dark-degeneracy mechanism explains both BAO shape blindness and SNe density sensitivity. The CMB result is conditional on the GR-calibrated compression and the adopted true-density ($\lambda=0$) transplant; varying $\lambda$ produces structural saturation within this mapping and makes fixed-map $\theta$ constraints protocol dependent. Profiling $\OL$ under the same protocol gives a descriptive summed per-probe minimum of $2.98$ near $\theta=5\deg$, with the adopted benchmark $0.14$ above it. On the finite benchmark branch the gradient-instability boundary lies at $z_*=7.877$ and obeys $(1+z_*)\propto\mFP^{2/3}$ in the self-similar regime. Moving it beyond recombination requires $\mFP/H_0\gtrsim4.1\times10^4$ in the prescribed mass-map cosmology. Every tabulated point of the scanned anchor set is scalar-unstable at $\mFP/H_0=25$ within the adopted linear/WKB diagnostic. A good background fit therefore does not, by itself, establish viability of the point or preference for the underlying theory. These statements are restricted to the tested benchmark and scanned domain; nonlinear saturation and observational viability remain undecided.
\end{abstract}
\medskip
